\documentclass{report}
\usepackage{amsmath,amssymb}
\setlength{\parindent}{0mm}
\setlength{\parskip}{1em}
\begin{document}
\begin{center}
\rule{6in}{1pt} \
{\large Alicia Klinvex \\
{\bf Clustering network data through effective use of eigensolvers and hypergraph models}}

Sandia National Labs \\ P O Box 5800 \\ Albuquerque \\ NM 87185-1320
\\
{\tt amklinv@sandia.gov}\\
Alicia Klinvex\\
Michael M. Wolf\\
Daniel Dunlavy\end{center}

In this talk, we discuss our efforts in leveraging eigensolvers (developed to
solve problems in computational science and engineering) in the spectral
analysis of graph and hypergraphs for community detection. In community
detection, the goal is to determine groupings of data objects given a set of
relationships amongst multiple objects. These relationships are often
represented by multiple edges in a graph in a simplified model, transforming the
relationship of many objects into multiple pairwise relationships between
objects. However, a multiway relationship between objects can be more
naturally represented by an hyperedge, a generalization of an edge that
contains one or more vertices. A focus of our work is understanding the
differences between these graph and hypergraph models.

One method of detecting communities is through spectral analysis of the
graph or hypergraph. A common procedure is to form the graph or
hypergraph Laplacian, compute a few of its eigenpairs, and run kmeans on
the subspace spanned by those eigenvectors. Let the graph incidence
matrix be $G$, the hypergraph incidence matrix be $H$, and the diagonal
matrices of vertex degrees and hyperedge cardinality be $D_v$ and $D_e$
respectively. Then the symmetric normalized graph Laplacian can be
expressed as
$L_G=I-D_{v_G}^{-1/2}\left(GG^T-D_{v_G}\right)D_{v_G}^{-1/2}$, and the
symmetric normalized hypergraph Laplacian is
$L_H=I-D_{v_H}^{-1/2}HD_e^{-1}H^TD_{v_H}^{-1/2}$. These operators may be
expressed explicitly as a matrix or implicitly as a series of linear algebra
operations. In practice, it may be unwise to explicitly build the matrix $L$,
especially with dynamic graphs that would require expensive updates to $L$
whenever the incidence matrix changed, so we have chosen to use the
implicitly defined operator in our code. We've implemented the above
methodology in both a Matlab prototype and (more recently) a high
performance computing software framework based implementation. In initial
experiments, we compared the spectral analysis resulting from the graph
and hypergraph models. For unweighted $G$, we found that the clustering
produced by performing spectral analysis on $L_G$ is considerably worse
than
that produced by performing spectral analysis on $L_H$; this is reflected both
in our Matlab results and in the existing literature. When the edges of
the graph corresponding to $G$ are weighted based on the cardinality of the
associated hyperedge, the graph results improve considerably and the
clustering quality is similar. However, this graph model still requires
significantly more storage, as well as more operations per matrix-vector
multiplication (for implicitly stored operators), so the hypergraph model
is computationally advantageous.

We chose to implement these spectral clustering methods in the Trilinos
framework, in order to solve these problems at a large scale (with an
eventual target of hundreds of millions to billions of vertices). We
leverage Trilinos's efficient parallel sparse eigensolver package
Anasazi, which contains many popular eigensolvers such as Block Krylov
Schur, LOBPCG, TraceMin, and the Riemannian Trust Region method. Anasazi
supports MPI+X, where X includes OpenMP, CUDA, and Pthreads, and it has
been used to solve data
analytics problems with over a billion vertices. We will be presenting results
obtained from our Trilinos-based code.

In recent work, we have been delving deeper into the eigensolver
framework, trying to better understand how to best use Anasazi in the
context of data science applications. One important issue is whether it
is more effective to compute the smallest eigenpairs of $L$ or the
largest eigenpairs of $I-L$. Although these problems have different
eigenvalues, the eigenvectors are the same and thus both can be used in
our spectral analysis. In general, it is easier to obtain the largest
eigenvalues of a matrix than it is to obtain the smallest ones, since
computing the smallest eigenpairs tends to require solving a series of
linear systems. However, $I-L_G$ may not be symmetric positive definite,
which can pose challenges for eigensolvers; $I-L_H$ has no such problem,
as it is guaranteed to be symmetric positive definite, which may be
advantagous to the hypergraph model. Another important issue that we will
discuss is how many eigenpairs should be provided to k-means and how
accurate they should be. The separation of the eigenvalues in the
spectrum is one contributing factor to the running time of the
eigensolver. Another factor is the tolerance requested from the
eigensolver. In this talk, we will explore these questions in the context
of Anasazi and the eigensolvers it provides, each of which have unique
pros and cons.


\end{document}
